\documentclass[11pt,a4paper]{article}

\usepackage[utf8]{inputenc}
\usepackage[T1]{fontenc}
\usepackage[lithuanian]{babel}
\usepackage{lmodern}
\usepackage{geometry}
\usepackage{xcolor}
\usepackage{enumitem}
\usepackage{titlesec}
\usepackage{fancyhdr}
\usepackage{microtype}
\usepackage{ragged2e}

\geometry{
  top=1.8cm,
  bottom=1.8cm,
  left=2.1cm,
  right=2.1cm
}

\definecolor{titleblue}{HTML}{1F3A5F}
\definecolor{ruleblue}{HTML}{6EA8E5}
\definecolor{callout}{HTML}{EEF2F6}

\setlength{\parindent}{0pt}
\setlength{\parskip}{0.42em}
\renewcommand{\familydefault}{\sfdefault}

\titleformat{\section}
  {\large\bfseries\color{titleblue}}
  {}{0pt}{}
\titlespacing*{\section}{0pt}{0.85em}{0.22em}

\setlist[itemize]{leftmargin=1.3em,itemsep=0.03em,topsep=0.12em,parsep=0pt}

\pagestyle{fancy}
\fancyhf{}
\fancyfoot[C]{\thepage}
\renewcommand{\headrulewidth}{0pt}
\renewcommand{\footrulewidth}{0pt}

\newcommand{\calloutbox}[1]{%
  \begin{center}
  \colorbox{callout}{%
    \parbox{0.94\textwidth}{\centering\bfseries #1}%
  }
  \end{center}
}

\begin{document}

\begin{center}
  {\LARGE\bfseries\color{titleblue} VU Šiaulių akademijos\\[2pt]
  Jaunųjų matematikų varžybos}\par
  \vspace{0.35em}
  {\color{ruleblue}\rule{0.92\textwidth}{0.8pt}}\par
  \vspace{0.45em}
  {\itshape Trumpas gidas pradedančiajam olimpiadininkui}
\end{center}

\section{Kas tai per varžybos?}

\textbf{VU Šiaulių akademijos Jaunųjų matematikų varžybos} -- tai tradicinis Lietuvos mokinių matematikos konkursas, skirtas \textbf{9--12 klasių mokiniams}.

Šios varžybos turi ilgą istoriją ir yra gerai žinomos tarp olimpiadine matematika besidominčių mokinių bei mokytojų.

Tai nėra paprastas mokyklinių žinių patikrinimas. Čia pateikiami \textbf{nestandartiniai, olimpiadinio pobūdžio uždaviniai}, kuriuose svarbu ne tik žinoti formules, bet ir gebėti ieškoti idėjos, pastebėti dėsningumus, argumentuoti ir kurti sprendimą.

Pradedančiam olimpiadininkui šios varžybos yra dar viena gera galimybė pasitikrinti savo gebėjimus rimtesnėje konkurencinėje aplinkoje.

\section{Kam jos skirtos?}

Varžybose dalyvauja \textbf{9--12 klasių mokiniai}.

Skirtingų klasių mokiniai varžosi savo amžiaus grupėse, todėl užduotys pritaikytos atitinkamam lygiui.

Dalyviai atvyksta iš įvairių Lietuvos miestų ir rajonų. Dažnai mokiniai į šias varžybas patenka per savo mokyklas, miestus ar rajonus, todėl dalyvavimas jose jau savaime reiškia, kad mokinys rimčiau domisi matematika.

\section{Kuo šios varžybos svarbios?}

Šios varžybos svarbios ne tik kaip atskiras matematikos konkursas.

Jos gali tapti \textbf{papildoma galimybe patekti į Lietuvos mokinių matematikos olimpiados šalies etapą}.

Tai ypač svarbu mokiniui, kuriam miesto etape nepavyko pasirodyti taip gerai, kaip jis tikėjosi.

Vienas nesėkmingas miesto etapas dar nereiškia, kad visas olimpiadinis sezonas baigtas.

Gerai pasirodžius Jaunųjų matematikų varžybose ir \textbf{užėmus prizinę vietą}, gali atsirasti \textbf{dar viena galimybė patekti į šalies etapą}.

\calloutbox{Miesto etape nepasisekė? Šios varžybos gali tapti antru šansu patekti į šalies etapą.}

\section{Kokie būna uždaviniai?}

Uždaviniai yra \textbf{olimpiadinio pobūdžio}.

Gali pasitaikyti:
\begin{itemize}
  \item algebros;
  \item geometrijos;
  \item skaičių teorijos;
  \item kombinatorikos;
  \item loginio ir nestandartinio mąstymo uždavinių.
\end{itemize}

Kaip ir kitose olimpiadinėse varžybose, svarbiausia nėra greitai pritaikyti išmoktą formulę.

\newpage

Dažniausiai reikia suprasti: \textbf{kokią pagrindinę šio uždavinio idėją?}

Kartais ji būna gana elementari, tačiau gerai paslėpta. Kitais atvejais reikia išbandyti kelis skirtingus kelius, kol pavyksta rasti tinkamą.

\section{Ar uždaviniai eina sunkėjimo tvarka?}

\textbf{Nebūtinai.}

Tai svarbus dalykas, kurį verta žinoti prieš varžybas.

Negalima manyti, kad pirmasis uždavinys būtinai bus lengviausias, o paskutinis -- sunkiausias.

Todėl gavus užduotis verta \textbf{pirmiausia perskaityti visą variantą}.

Gali būti, kad vienas iš vėlesnių uždavinių tau bus daug aiškesnis negu pirmasis.

Šiuo požiūriu Jaunųjų matematikų varžybos panašesnės į Matulionio konkursą nei į Vilniaus miesto olimpiadą, kur uždavinių tvarka dažniau padeda orientuotis jų sunkume.

\section{Kaip vertinami sprendimai?}

Svarbus ne tik galutinis atsakymas, bet ir \textbf{sprendimo pagrindimas}.

Reikia parodyti, kodėl tavo teiginiai yra teisingi ir kaip priėjai prie galutinės išvados.

Todėl verta stengtis rašyti sprendimą taip, kad jį galėtų perskaityti ir suprasti kitas matematikas.

Jeigu viso uždavinio išspręsti nepavyksta, vis tiek verta užrašyti prasmingą pažangą:
\begin{itemize}
  \item svarbią idėją;
  \item įrodytą tarpinį teiginį;
  \item rastą dėsningumą;
  \item išnagrinėtą atvejį;
  \item teisingai gautą tarpinį rezultatą.
\end{itemize}

Olimpiadinėje matematikoje \textbf{dalinis sprendimas taip pat gali būti vertingas}.

\section{Ką svarbiausia prisiminti pradedančiajam?}

Šios varžybos nėra tik dar vienas konkursas kalendoriuje.

Jos yra gera proga:
\begin{itemize}
  \item pasitikrinti save;
  \item gauti daugiau rimtų olimpiadinių uždavinių patirties;
  \item palyginti savo lygį su kitais stipriais mokiniais;
  \item išmokti geriau valdyti laiką;
  \item išmokti aiškiau rašyti sprendimus.
\end{itemize}

Ir dar svarbiau -- jos gali suteikti \textbf{dar vieną galimybę patekti į šalies etapą}, jeigu miesto olimpiadoje nepavyko pasirodyti taip gerai, kaip norėjosi.

Todėl po nesėkmingo miesto etapo nereikėtų manyti, kad viskas baigta.

Olimpiadinėje matematikoje labai svarbu mokėti ne tik laimėti, bet ir \textbf{grįžti po nesėkmės, toliau spręsti ir pasinaudoti kita galimybe}.

\begin{center}
\textbf{Jaunųjų matematikų varžybos yra būtent viena iš tokių galimybių.}
\end{center}

\end{document}
